\chapter{Cli}
\section{Tổng quan Giao diện Dòng lệnh (CLI)}
Quản trị

`KBMS.CLI` là một ứng dụng Console (C\#) mạnh mẽ, cung cấp khả năng tương tác trực tiếp với máy chủ KBMS thông qua giao thức nhị phân và dòng lệnh KQL/KDL.

\subsection{Các Tính năng Chính}

*   \textbf{REPL (Read-Eval-Print Loop)}: Người dùng nhập câu lệnh, CLI gửi tới Server, nhận kết quả và in ra ngay lập tức.
*   \textbf{Advanced Line Editing}: Hỗ trợ đầy đủ các phím di chuyển và thao tác nhanh (`LineEditor.cs`):
    *   \textbf{Mũi tên Lên/Xuống}: Duyệt lịch sử lệnh.
    *   \textbf{Home/End}: Di chuyển nhanh đến đầu/cuối dòng.
    *   \textbf{Delete/Backspace}: Xóa ký tự tại con trỏ.
    *   \textbf{Esc x2}: Xóa sạch bộ đệm đang nhập (Double Escape to Clear).
*   \textbf{Multi-line Input}: Hỗ trợ nhập các câu lệnh dài. Chế độ thụt đầu dòng tự động `->` giúp phân biệt dòng tiếp nối.
*   \textbf{Display Modes}: Cung cấp hai chế độ hiển thị chính thông qua `ResponseParser.cs`:
    *   \textbf{Table Mode}: Hiển thị bảng ngang (MySQL style).
    \textit{   }*Vertical Mode (\G)\textit{}: Hiển thị theo hàng dọc, tự động kích hoạt cho các lệnh `DESCRIBE` hoặc `EXPLAIN` để tối ưu hóa việc đọc cấu trúc tri thức phức tạp.
*   \textbf{Chế độ Chạy Script (`SOURCE`)}: Cho phép thực thi các tệp tin truy vấn tri thức lớn (`.kbql`) một cách tự động, tự dừng khi gặp lỗi.
*   \textbf{Xử lý Phản hồi Phức tạp}: Sử dụng `ResponseParser.cs` để hiển thị kết quả truy vấn dưới dạng bảng (Table), đồ thị tri thức hoặc thông báo lỗi line-accurate.
*   \textbf{Tự động kết nối lại (Auto-reconnect)}: CLI tự động thử kết nối lại nếu máy chủ bị ngắt kết nối đột ngột (Heartbeat timeout).

---

\subsection{Các Lệnh Hệ thống Đặc biệt}

| Lệnh | Mô tả |
| :--- | :--- |
| `LOGIN <u|p>` | Đăng nhập bảo mật (mật khẩu không lưu vào lịch sử). |
| `SOURCE <path>` | Thực thi script từ tệp tin bên ngoài. |
| `CONNECT` | Kết nối lại máy chủ thủ công. |
| `CLEAR` | Xóa sạch màn hình console. |

---

\subsection{Trải nghiệm người dùng (UX)}

Để mang lại trải nghiệm chuyên nghiệp như một hệ quản trị CSDL thực thụ, CLI được tích hợp các phím tắt:
*   \textbf{Mũi tên Lên/Xuống}: Xem lại lịch sử lệnh đã gõ.
*   \textbf{Home/End}: Di chuyển con trỏ nhanh trong dòng lệnh.
*   \textbf{Esc x2}: Xóa nhanh input đang nhập dở.

> [!TIP]
> \textbf{LineEditor \& HistoryManager}
> Hệ thống quản lý lịch sử lệnh của KBMS CLI được thiết kế để bỏ qua các lệnh nhạy cảm như `LOGIN` hoặc `CREATE USER`, đảm bảo an toàn thông tin tối đa cho người dùng. 🛡️

\section{Luồng Xử lý Dữ liệu CLI}

Để đảm bảo hiệu năng trong môi trường dòng lệnh, `KBMS.CLI` vận hành theo một luồng xử lý phi đồng bộ (Asynchronous) dựa trên socket.

\subsection{Luồng Gửi lệnh (Command Execution Flow)}

Khi người dùng nhấn Enter, CLI thực hiện các bước sau:

![cli\_input\_flow.mmd](../assets/diagrams/cli\_input\_flow.png)

1.  \textbf{Input Buffering}: Tích lũy các dòng cho đến khi gặp dấu `;`.
2.  \textbf{Request Construction}: Đóng gói lệnh thành các `Message` (Dạng `QUERY` hoặc `LOGIN`).
3.  \textbf{Binary Transport}: Gửi qua `KBMS.Network` (Socket).
4.  \textbf{Wait for Response}: Lắng nghe phản hồi từ máy chủ.

---

\subsection{Bộ máy Phân tích Phản hồi (`ResponseParser`)}

Phần phức tạp nhất của CLI nằm ở `ResponseParser.cs`. Do kết quả từ Server có thể là một luồng dữ liệu (Streaming ROWS), CLI phải xử lý từng gói tin:

*   \textbf{METADATA}: Định nghĩa các cột (Name, Type).
*   \textbf{ROW}: Dữ liệu thực tế cho từng bản ghi.
*   \textbf{RESULT}: Thông điệp cuối cùng (v.d: "Insert successful").
*   \textbf{ERROR}: Chứa thông tin lỗi (Message, Line, Column).

\subsubsection{Quy trình hiển thị:}
`ResponseParser.cs` thực hiện vẽ bảng "động" theo thuật toán tối ưu:
1.  \textbf{Header Rendering}: Ngay khi nhận `METADATA`, CLI tính toán độ rộng cột lớn nhất dựa trên tên cột và vẽ khung Header.
2.  \textbf{Multi-line Cell Support}: Nếu dữ liệu trong một ô chứa dấu xuống dòng (`\n`), CLI sẽ tự động chia ô đó thành nhiều dòng và vẽ đường kẻ phân cách hàng (`|---+---|`) để đảm bảo bảng luôn cân đối.
3.  \textbf{Vertical Rendering}: Đối với các phản hồi thuộc nhóm `Explain\_` hoặc `Describe\_`, CLI chuyển sang chế độ hiển thị theo cặp `Field: Value` trên từng hàng dọc.

---

\subsection{Chế độ Chạy Script (Batch Source)}

CLI hỗ trợ thực thi khối lượng lớn lệnh thông qua tệp tin `.kbql`. Luồng này được xử lý tuần tự để đảm bảo tính nhất quán:

![uc\_cli\_batch\_source.png](../assets/diagrams/uc\_cli\_batch\_source.png)

---

\subsection{Đa luồng & Kết nối (Multi-threading)}

CLI chạy hai luồng chính để duy trì trạng thái:
1.  \textbf{Main Thread}: Chờ nhận input của người dùng.
2.  \textbf{Heartbeat Thread}: (Tùy chọn) Gửi tín hiệu duy trì kết nối nếu cần thiết.

> [!IMPORTANT]
> \textbf{Streaming Mode}
> Nhờ cơ chế nhận tin phi tập trung, người dùng có thể thấy kết quả trả về ngay lập tức cho các truy vấn lớn (`LIMIT 10000`) mà không cần chờ Server nén toàn bộ dữ liệu. ⚡

\section{Các Trường hợp Sử dụng CLI}

KBMS CLI là công cụ tối thượng cho các quản trị viên hệ thống và chuyên gia tri thức cần sự tốc độ và khả năng tự động hóa.

\subsection{Kịch bản 1: Phân tích Cấu trúc (Structure Analysis)}

*   \textbf{Mục tiêu}: Xem chi tiết định nghĩa của một Concept phức tạp với hàng chục biến.
*   \textbf{Sơ đồ Hoạt động}:

![uc\_cli\_vertical\_mode.png](../assets/diagrams/uc\_cli\_vertical\_mode.png)

*   \textbf{Quy trình chi tiết}:
    1.  \textbf{Gửi lệnh}: Người dùng gõ `DESCRIBE <ConceptName>;`.
    2.  \textbf{Nhận diện}: CLI nhận gói tin Metadata có thuộc tính `ConceptName` bắt đầu bằng `Describe\_`.
    3.  \textbf{Chế độ Dọc}: `ResponseParser` tự động chuyển sang Vertical Mode để hiển thị danh sách biến theo chiều dọc, tránh việc bảng bị tràn màn hình (Overflow).
    4.  \textbf{Kết quả}: Từng biến được liệt kê rõ ràng kèm kiểu dữ liệu tương ứng.

\subsection{Kịch bản 2: Tự động hóa Nạp Tri thức (Batch Loading)}

*   \textbf{Mục tiêu}: Nạp hàng ngàn thực thể tri thức từ tệp tin script có sẵn.
*   \textbf{Quy trình chi tiết}:
    1.  \textbf{Chuẩn bị}: Tạo tệp `import.kbql` chứa các lệnh `INSERT`.
    2.  \textbf{Thực thi}: Gõ `SOURCE import.kbql;` tại CLI.
    3.  \textbf{Vòng lặp}: CLI đọc từng block lệnh (phân tách bởi dấu `;`) và gửi tới Server.
    4.  \textbf{Xử lý sai sót}: Nếu một lệnh gặp lỗi, CLI sẽ dừng ngay lập tức và chỉ ra vị trí dòng lệnh bị lỗi trong tệp nguồn để quản trị viên sửa chữa.
    5.  \textbf{Thống kê}: Sau khi hoàn thành, CLI in ra tổng số bản ghi đã nạp và tổng thời gian thực thi.

---

\subsection{Kịch bản 3: Kiểm tra Kết nối & An ninh (Security Check)}

*   \textbf{Mục tiêu}: Đăng nhập và kiểm tra tình trạng kết nối tới máy chủ từ xa.
*   \textbf{Quy trình}:
    1.  Người dùng gõ `LOGIN <username> <password>`.
    2.  CLI gửi gói tin `LOGIN` bảo mật tới Server.
    3.  Server xác thực và trả về `LOGIN\_SUCCESS`.
    4.  Prompt của CLI chuyển từ `login>` sang `kbms> `, sẵn sàng cho các lệnh truy vấn.

> [!TIP]
> \textbf{Hiệu năng Dòng lệnh}
> Nhờ cơ chế streaming dữ liệu, CLI có thể hiển thị những bản ghi đầu tiên của một truy vấn `SELECT` lớn chỉ sau vài mili giây, giúp người dùng đánh giá dữ liệu mà không cần chờ đợi toàn bộ kết quả. ⚡

\section{Xác thực CLI (CLI Validation)}

KBMS CLI là giao diện dòng lệnh mạnh mẽ, cho phép tương tác trực tiếp và chạy các kịch bản tri thức quy mô lớn.

\subsection{Kiểm thử Phiên làm việc (Session Logs)}

Mọi thao tác trong CLI đều được ghi nhận và xác thực thông qua log kết nối tới Server.

*   \textbf{Handshake nhị phân}: CLI gửi bản tin định danh và nhận SESSION\_ID.
*   \textbf{Batch execution}: Chạy script qua lệnh `SOURCE`.

\subsubsection{Minh chứng Mã nguồn (CLI System):}
\textbf{[Missing Image: code_test_cli.png]}\\ \textit{Hình 12.1: Minh chứng mã nguồn kiểm thử giao diện dòng lệnh (CLI).}

\subsubsection{Minh chứng Kết quả (Session Log):}
\textbf{[Missing Image: terminal_test_cli_query.png]}\\ \textit{Hình 12.1: Giao diện tương tác CLI thực thi câu lệnh truy vấn tri thức.}

\subsection{Kiểm thử Định dạng Kết quả}

Xác thực khả năng căn chỉnh bảng và hiển thị kết quả truy vấn dưới dạng trực quan.

| Query | Kết quả mô phỏng | Hiển thị bảng |
| :--- | :--- | :--- |
| `SELECT *` | 3 Rows | Bảng có header lam (Cyan) |
| `DESCRIBE` | Metadata | Bảng thông tin Concept |

\textbf{[Missing Image: placeholder_cli_describe_output.png]}\\ \textit{Hình 10.3: Hiển thị kết quả truy vấn dưới dạng bảng (Table View) trong CLI.}

---

> [!NOTE]
> Bạn có thể chạy `full\_test.kbql` trực tiếp từ CLI để kiểm chứng toàn bộ 96 kịch bản thử nghiệm chỉ trong tích tắc.

